\section{Introduction}
\label{sec:intro}

\subsection{From encoder robustness to infrastructure-level adaptation}

Dense retrieval has quietly become the backbone of a much larger system than
a search box: it is the retrieval layer of retrieval-augmented generation
(RAG) pipelines that ground large language models in external text. In this
setting, three pressures are converging in a way that the corruption-robustness
literature has treated mostly in isolation.

First, the encoder itself is increasingly \emph{not owned} by the team
building the retrieval system. Production RAG stacks routinely embed
documents and queries with a third-party API (a commercial embedding
endpoint or a large shared open-weight encoder) that cannot be fine-tuned,
distilled, or even inspected beyond its output vectors. Robustness and
compression techniques that assume gradient access to the encoder, or that
require re-embedding the corpus under a new objective, are not deployable in
this regime. Whatever adaptation is applied must live \emph{after} the
frozen encoder, as a lightweight map on top of embeddings that already exist.

Second, the queries reaching these systems are no longer curated search
strings. They arrive from voice transcription, OCR'd screenshots, multilingual
users typing in a non-native script, and increasingly from other language
models paraphrasing or rewriting a user's intent before retrieval is even
called. The corruption statistics of this traffic are corpus- and
deployment-specific: they are not captured by a fixed, published perturbation
benchmark, and they change as the upstream product changes. A retrieval
adaptation layer that must be re-trained by the encoder's original authors,
on the original training corpus, cannot track this drift; one that can be
re-estimated by the deploying team from their own paired clean/noisy traffic
can.

Third, the storage and latency budget of RAG systems scales with corpus size
far more aggressively than the original ad-hoc search literature anticipated:
production vector stores now index tens to hundreds of millions of
passages, and every added dimension multiplies index memory, network
transfer, and approximate nearest-neighbor search cost. Compression is no
longer a secondary concern relative to accuracy; it is a load-bearing
requirement, and any accuracy loss it introduces is compounded exactly when
the traffic is already degraded by corruption.

These three pressures point to the same design target: a small,
cheaply re-estimated, encoder-agnostic linear layer that sits between a
frozen embedding model and a nearest-neighbor index, compresses the vectors
it is given, and is explicitly fit to the gap between an organization's own
clean documents and the corrupted queries its own users actually produce.
This is a narrower and more operational target than ``robust
representation learning'' in general: it does not ask for a better encoder,
a better loss to train one, or a universal corruption benchmark. It asks
what can be done \emph{post-encoder, pre-index}, with only paired clean and
corrupted embeddings and no labels, no queries, and no access to model
weights.

\subsection{Concrete deployment scenarios}
\label{sec:scenarios}

Three recurring production scenarios make the target concrete and
distinguish it from a purely academic corruption benchmark.

\emph{Scenario A: multilingual customer-support RAG.} A support assistant
embeds a knowledge base once with a hosted embedding API and serves queries
typed by non-native speakers, frequently containing transliteration
artifacts and autocorrect substitutions. The organization has no access to
the encoder's weights and no ability to re-embed the knowledge base under a
different objective, but it does have logs pairing clean knowledge-base
passages with the corrupted phrasings that its own users actually type.
This is precisely the paired, label-free, query-free fitting regime the
method requires.

\emph{Scenario B: voice-driven retrieval.} An assistant that accepts speech
input routes transcripts through an automatic-speech-recognition system
before retrieval. ASR error patterns are systematic (homophones, digit and
proper-noun substitution) rather than uniform random noise, so a fixed
published corruption suite is a poor proxy; the deploying team can, however,
synthesize or log paired clean/ASR-corrupted embeddings from its own
domain and re-estimate the adaptation layer as that error profile evolves
with the ASR model.

\emph{Scenario C: agentic query rewriting.} Increasingly, the text that
reaches an embedding call is itself the output of another language model
that has paraphrased or expanded a user's request. This introduces a
distinct, model-dependent corruption process at the embedding boundary
that changes whenever the upstream rewriting model changes, again favoring
an adaptation layer that a downstream team can refit cheaply rather than
one baked into the encoder at pretraining time.

Across all three, the operative constraints are the same: no relevance
labels, no access to evaluation queries during fitting, no gradient access
to the encoder, and a corruption process that is local to the deployment
rather than universal. These constraints, rather than a specific corruption
taxonomy, are what should be read as defining the problem class this paper
addresses; the OCR-like and query-variant corruptions used in the reported
experiments are one sampling of that class, not its definition.

\subsection{Why a metric-learning view, not a denoising view}

A natural first instinct is to treat this as a denoising problem: learn a
map that pulls a corrupted embedding back toward its clean counterpart and
judge it by reconstruction error. That framing, however, is misaligned with
how the adapted embeddings are actually consumed. A RAG retrieval layer
never inspects a reconstructed vector directly; it ranks documents by a
similarity score computed from projected vectors. Two linear maps with
identical reconstruction subspaces can induce different rankings, and a map
that reconstructs well on held-out corruption can still degrade ranking
if it does so along directions the similarity score weights differently
than the loss did. Compression, corruption-adjustment, and ranking
therefore need to be posed as a single object: a low-rank, positive
semidefinite similarity structure, together with an explicit account of how
that structure relates to (i) the corpus covariance geometry used to select
it, (ii) the loss geometry used to refine it, and (iii) the geometry
actually queried at serving time. Separating these three notions of
geometry, rather than conflating them under a single ``robust embedding''
label, is what allows the adaptation layer to be reasoned about and
diagnosed rather than only benchmarked.

\subsection{Positioning relative to adjacent lines of work}

\paragraph{Encoder fine-tuning and contrastive retraining.} Directly
fine-tuning or contrastively re-training the encoder on noisy queries
gives the strongest ceiling on robustness when it is available, but is
foreclosed exactly in the black-box-encoder regime described above, and
must be repeated whenever the corruption distribution drifts. The method
studied here targets the complementary, encoder-frozen regime.

\paragraph{Unsupervised dimensionality reduction.} Classical compression
such as PCA or random projection is encoder-agnostic and cheap to
re-estimate, but is fit to corpus variance alone; it has no mechanism for
down-weighting directions that a specific corruption process happens to
inflate, because it never sees paired corrupted views.

\paragraph{Query-side robustness heuristics.} Spelling correction,
query rewriting, and retrieval-time ensembling address corruption at the
query in isolation, at inference latency cost, and do not touch the
compression problem at all; they are compatible with, rather than a
substitute for, an adapted shared embedding space for documents and
queries alike.

\paragraph{Classical noise-adjusted spectral estimators.} The generalized
eigenvalue problem underlying signal-to-noise-ratio maximization is decades
old and is not itself a contribution of this work; what is missing from
that lineage is a treatment of the paired \emph{textual} corruption setting,
the retrieval-specific low-rank cosine metric it induces, and an account
of when subsequent geometric refinement helps or actively harms the
metric it starts from.

\paragraph{Embedding compression and quantization.} Product quantization,
scalar quantization, and nested (``Matryoshka''-style) representations
reduce index size directly and are attractive precisely because they are
also encoder-agnostic. They compress by discarding precision or by
truncating a representation the encoder was trained to keep informative at
every prefix length; they are not fit to any paired corruption signal and
so provide no mechanism for asking which retained directions are
noise-bearing for a specific deployment. The adaptation layer studied here
is complementary rather than competing: its output is itself a
lower-dimensional real vector that remains a candidate for downstream
quantization.

\paragraph{Multi-vector and cross-encoder reranking.} Late-interaction
multi-vector retrievers and cross-encoder rerankers improve robustness and
accuracy by spending more compute per query at serving time, typically
after an initial single-vector retrieval pass. This is a serving-time
robustness strategy rather than an embedding-adaptation strategy, incurs
materially higher query latency and index size, and is naturally staged
\emph{after} an adapted first-pass retriever rather than as a replacement
for one.

\paragraph{Adversarial and distributionally robust training.} Robust
optimization formulations that minimize worst-case loss over a corruption
neighborhood again presuppose access to the encoder's parameters and a
specified perturbation set at training time; they target the same
robustness goal from the retrainable-encoder side of the divide this paper
is explicitly on the other side of.

\subsection{Contributions}
\label{sec:contributions}

Framed against the infrastructure target above, the contributions of the
paper are the following.

\begin{enumerate}[leftmargin=1.6em]
 \item \textbf{A three-geometry decomposition of a post-hoc adaptation
 layer.} We separate the Euclidean reconstruction projector, the
 corruption-constraint geometry fixed by the estimated noise covariance,
 and the low-rank cosine metric actually consulted at serving time, and
 show that they answer three different questions rather than being
 interchangeable descriptions of the same object. This decomposition is
 what lets an operator diagnose \emph{which} stage of a deployed adaptation
 layer is responsible for a given behavior, rather than treating the layer
 as a single opaque compressor.

 \item \textbf{A basis-sensitivity result with a direct operational
 consequence.} We show precisely which changes of basis leave the
 reconstruction subspace unchanged but alter every reduced-space cosine
 score, and identify the conformal condition under which ranking is
 basis-invariant. The practical corollary is a concrete, checkable
 explanation for why a standard orthonormalization step (Euclidean QR),
 applied for numerical convenience, can silently degrade a deployed
 retriever even though it preserves the span an engineer might think of as
 ``the same subspace.''

 \item \textbf{An exact statistical account of the compression--robustness
 trade-off.} The denoising-risk decomposition gives the two population
 quantities -- discarded clean variance and transmitted noise variance --
 that any rank-$r$ projector trades against each other, giving deployment
 teams a named quantity to reason about (and eventually estimate from their
 own logs) instead of only a benchmark accuracy number.

 \item \textbf{A local geometric account of what the training objective is
 actually doing.} The damped pullback metric shows how the reconstruction,
 invariance, and noise-suppression terms jointly define a single local
 notion of distance on the constraint manifold, with an explicit variance
 identity showing that the contrastive term concentrates on relative
 rather than absolute score changes and that its curvature is
 temperature-dependent. This clarifies what is being optimized without
 requiring the reader to treat the composite loss as an unexamined sum of
 heuristics.

 \item \textbf{A margin-based sufficient condition for ranking stability}
 that connects the noise-suppression term directly to a top-1 stability
 guarantee under bounded embedding perturbation, giving a mechanistic
 rather than purely correlational explanation for why suppressing
 projected noise energy should help ranking under corruption.

 \item \textbf{A finite-sample estimation-error bound that separates the
 two estimation problems an operator actually faces} -- clean-covariance
 estimation and corruption-covariance estimation -- rather than bounding
 their combined effect, so that the sample budget for clean documents and
 the sample budget for paired corrupted traffic can in principle be
 reasoned about separately.

 \item \textbf{An explicit statement of evidentiary scope.} Each of the
 above results is paired with an account of exactly what the reported,
 single-seed, single-encoder experiments do and do not establish about it,
 and what a deployment-oriented team would additionally need to check
 (multiple seeds, multiple encoders, out-of-distribution corruption,
 sample-size sensitivity) before relying on the method in production.

 \item \textbf{Infrastructure-oriented theoretical extensions beyond the
 core estimator.} Motivated directly by the deployment scenarios of
 Section~\ref{sec:scenarios}, Section~\ref{sec:extensions} develops four
 additional results that the corruption-robustness framing alone does not
 call for: an explicit, non-asymptotic sample-complexity rate for the two
 covariance-estimation problems underlying Theorem~\ref{thm:finite}; a
 cross-deployment transfer bound quantifying the excess retrieval risk of
 serving an adaptation layer fit on one corruption distribution against a
 different production corruption distribution; a sequential monitoring
 test, built from the same concentration machinery, that flags when a
 deployed layer should be re-fit without requiring relevance labels or
 evaluation queries; and a label-free rank-selection guarantee for the
 isotropic-noise case, with an explicit account of why the anisotropic case
 requires the corruption-metric projector rather than the Euclidean one.
\end{enumerate}

\subsection{Forward-looking framing and scope of the contribution}

The remainder of the paper is organized around this infrastructure-level
target rather than around corruption robustness as an end in itself.
Section~\ref{sec:problem} fixes the frozen-encoder, paired-view setting in
which the adaptation layer must operate without queries or labels.
The methods section develops the adaptation layer as three deliberately
distinguished geometric objects -- a Euclidean reconstruction projector, a
corruption-constraint geometry, and the low-rank cosine metric actually
used at serving time -- together with a closed-form initializer and an
optional geometry-aware refinement stage compatible with either Euclidean
or corruption-metric retraction. The analysis section then asks, for each
of these pieces, exactly what can and cannot be claimed about it prior to
deployment: which optimum is global, which trade-off is an exact identity
versus an empirical trend, which retraction choice is guaranteed to
preserve which invariant, and what finite-sample conditions would be needed
for the estimated subspace to be close to its population counterpart. This
separation between architectural claims and evidentiary status is intended
to make the method auditable as a piece of retrieval infrastructure -- one
that a team operating a black-box-encoder RAG system could re-fit on its
own traffic, reason about, and monitor -- rather than as a single
benchmark-driven robustness trick tied to one encoder and one published
corruption suite.

This auditability goal is deliberately treated as a first-class design
criterion rather than as an afterthought to a benchmark result. As retrieval
adaptation layers move from research artifacts into components that sit
inside production language-model systems, it becomes increasingly relevant
to be able to say, for a given deployed layer, which of its guarantees are
exact algebraic facts (feasibility of the retraction, global optimality of
the closed-form stage, positive-definiteness of the local task metric),
which are statistical identities that hold under stated assumptions (the
denoising-risk decomposition), and which are still empirical trends awaiting
further confirmation (the benefit of geometric refinement over the
closed-form stage alone, the practical value of task-metric preconditioning
as an explicit optimizer, and generalization across encoders and corruption
families beyond the mismatch regime analyzed below).

Two of the three directions raised by this stance -- what happens when the
corruption distribution used to fit the layer differs from the one it is
served on, and how a deployed layer can be monitored for exactly this kind
of drift without relevance labels -- are not left as unformalized future
work here. Section~\ref{sec:extensions} develops them as explicit
theoretical extensions: a finite-sample rate that turns the abstract
estimation errors of Theorem~\ref{thm:finite} into concrete sample-size
requirements, a cross-deployment transfer bound that quantifies the excess
risk of serving a layer fit on one corruption distribution against a
different deployment distribution, a sequential statistical test that uses
that same concentration machinery to flag when re-fitting is warranted, and
a label-free rank-selection rule for the isotropic-noise case. The remaining
open direction -- using the damped pullback metric of
Theorem~\ref{thm:taskmetric} as an explicit natural-gradient preconditioner,
and validating all of the above jointly across multiple frozen encoders and
corruption families -- is left as future work in the sense described above:
a concrete, falsifiable target rather than an unformalized aspiration.
Positioning the contribution this way is what makes it a piece of retrieval
infrastructure for the RAG era rather than a corruption-benchmark result in
isolation.

\section{Infrastructure-Oriented Theoretical Extensions}
\label{sec:extensions}

The results in Section~\ref{sec:analysis} analyze the estimator, projector,
retraction, and composite objective as fit on a single, fixed dataset. The
deployment scenarios of Section~\ref{sec:scenarios}, however, raise three
questions that a single-fit analysis does not address: how many paired
samples are actually required before the finite-sample guarantee of
Theorem~\ref{thm:finite} is meaningful; what happens to retrieval risk when
the corruption distribution used to fit the adaptation layer differs from
the one it is later served on; and how a deployed layer can be told, from
its own traffic and without relevance labels, that it has drifted far
enough from its fitting distribution to warrant re-estimation. This section
answers all three, together with a fourth, narrower question -- when the
rank $r$ itself can be chosen without labels -- and states plainly which
parts of the answer are exact and which remain conditional on assumptions
that a deploying team would need to check against its own traffic.

\subsection{Setting and additional notation}
\label{sec:extensions-setting}

We retain the notation of Section~\ref{sec:analysis}: clean embeddings have
covariance $\Sigma_s$, paired corruption residuals have covariance
$\Sigma_n$, and $B\succ0$ is a regularized version of $\Sigma_n$ as in
\cref{eq:B}. Throughout this section we additionally assume:
\begin{itemize}[leftmargin=1.6em]
 \item[(A1)] Clean embedding rows $x_i-\mu_s$, $i=1,\dots,N$, are
 independent, centered, sub-Gaussian random vectors in $\R^d$ with
 covariance $\Sigma_s$ and sub-Gaussian (Orlicz) norm at most $K_s$.
 \item[(A2)] Paired residuals $e_i^{(v)}-\mu_n$, $i=1,\dots,N$,
 $v=1,\dots,V$, are independent, centered, sub-Gaussian random vectors in
 $\R^d$ with covariance $\Sigma_n$ and sub-Gaussian norm at most $K_n$,
 independent of the clean embeddings.
 \item[(A3)] $B\succeq\beta\I_d$ and $\widehat B\succeq\beta\I_d$ for a
 known or estimable lower bound $\beta>0$, as already assumed in
 Theorem~\ref{thm:finite}.
\end{itemize}
Assumptions (A1)--(A2) are standard in high-dimensional covariance
estimation \citep{vershynin2018high} and are satisfied, for instance, when
embeddings are bounded (as is the case after any bounded pooling or
normalization layer in a text encoder) or approximately Gaussian, which is
a common empirical observation for pretrained sentence embeddings even
though it is not guaranteed by construction. We write $n=NV$ for the
number of paired-residual samples and reserve $N$ for the number of clean
samples, matching the estimators of \cref{eq:B}.

\subsection{Sample complexity for the corruption-adjusted subspace}
\label{sec:sample-complexity}

Theorem~\ref{thm:finite} states its guarantee in terms of abstract operator
norm errors $\epsilon_c=\norm{\widehat\Sigma_s-\Sigma_s}_2$ and
$\epsilon_B=\norm{\widehat B-B}_2$. It does not say how many samples are
needed to make these errors small. The following theorem closes this gap
under (A1)--(A2).

\begin{theorem}[Non-asymptotic covariance concentration]
\label{thm:sampleconcentration}
Under (A1)--(A2), there exists an absolute constant $C>0$ such that for
every $t\ge0$, with probability at least $1-2e^{-t}$,
\begin{equation}
 \epsilon_c=\norm{\widehat\Sigma_s-\Sigma_s}_2
 \;\le\;
 CK_s^2\norm{\Sigma_s}_2
 \left(\sqrt{\frac{d+t}{N}}+\frac{d+t}{N}\right),
 \label{eq:concentration-clean}
\end{equation}
and, independently, with probability at least $1-2e^{-t}$,
\begin{equation}
 \norm{\widehat\Sigma_n-\Sigma_n}_2
 \;\le\;
 CK_n^2\norm{\Sigma_n}_2
 \left(\sqrt{\frac{d+t}{n}}+\frac{d+t}{n}\right).
 \label{eq:concentration-noise}
\end{equation}
Moreover, the shrinkage-and-ridge map $\widehat\Sigma_n\mapsto B=
\mathcal S_{\alpha_n}(\widehat\Sigma_n)+\rho\I_d$ of \cref{eq:B} is
$1$-Lipschitz in operator norm, so \cref{eq:concentration-noise} also
bounds $\epsilon_B=\norm{\widehat B-B}_2$, with $B$ interpreted as the
population quantity obtained by applying the same shrinkage-and-ridge map
to $\Sigma_n$.
\end{theorem}
\begin{proof}
Inequalities \eqref{eq:concentration-clean}--\eqref{eq:concentration-noise}
are the standard sub-Gaussian sample-covariance concentration bound
\citep[Theorem~4.7.1]{vershynin2018high}, applied once to the $N$ i.i.d.\
clean samples and once to the $n=NV$ i.i.d.\ paired residuals; independence
across the two applications follows from the independence assumed in
(A1)--(A2). For the Lipschitz claim, write $\Delta=\widehat\Sigma_n-
\Sigma_n$. By the triangle inequality applied twice,
\begin{align*}
 \norm{\mathcal S_{\alpha_n}(\widehat\Sigma_n)-\mathcal S_{\alpha_n}(\Sigma_n)}_2
 &=\Big\|(1-\alpha_n)\Delta
 +\alpha_n\frac{\Tr(\Delta)}{d}\I_d\Big\|_2\\
 &\le(1-\alpha_n)\norm{\Delta}_2+\alpha_n\frac{|\Tr(\Delta)|}{d}
 \;\le\;(1-\alpha_n)\norm{\Delta}_2+\alpha_n\norm{\Delta}_2
 =\norm{\Delta}_2,
\end{align*}
using $|\Tr(\Delta)|\le d\norm{\Delta}_2$. Adding the same deterministic
ridge $\rho\I_d$ to both the empirical and population quantities leaves
their difference, and hence its operator norm, unchanged.
\end{proof}

\begin{corollary}[Sample size for a target subspace accuracy]
\label{cor:samplesize}
Fix a target accuracy $\eta\in(0,1)$ and failure probability $\delta\in
(0,1)$, and suppose the population eigengap $g=\lambda_r(C)-\lambda_{r+1}(C)$
of Theorem~\ref{thm:finite} is strictly positive. If
\begin{equation}
 N\;\ge\;
 c_1K_s^2\frac{\norm{\Sigma_s}_2^2}{\beta^2\eta^2g^2}
 \left(d+\log\tfrac{4}{\delta}\right),
 \qquad
 NV\;\ge\;
 c_2K_n^2\frac{\norm{\Sigma_n}_2^2\norm{\Sigma_s}_2^2}
 {\beta^4\eta^2g^2}
 \left(d+\log\tfrac{4}{\delta}\right)
 \label{eq:samplesizereq}
\end{equation}
for absolute constants $c_1,c_2>0$, then, with probability at least
$1-\delta$, the right-hand side of \cref{eq:whitenedperturb} is at most
$g/2$ and
\begin{equation}
 \norm{\sin\Theta(\widehat V_r,V_r)}_2\;\le\;\eta.
\end{equation}
\end{corollary}
\begin{proof}
Apply Theorem~\ref{thm:sampleconcentration} with $t=\log(4/\delta)$ to both
$\epsilon_c$ and $\epsilon_B$ and a union bound over the two events, each
of probability at least $1-\tfrac{\delta}{2}$. Substituting the resulting
bounds into \cref{eq:whitenedperturb} and solving for $N$ and $NV$ so that
$\epsilon_c/\beta\le g\eta/4$ and $\norm{\Sigma_s}_2\epsilon_B/\beta^2\le
g\eta/4$ (so that their sum is at most $g\eta/2\le g/2$ whenever $\eta\le1$,
and is at most $g\eta$ after the final rescaling of the finite-sample bound
\cref{eq:finitebound}) yields \cref{eq:samplesizereq} up to the absolute
constants $c_1,c_2$; \cref{eq:finitebound} then gives the stated subspace
bound.
\end{proof}

\paragraph{Connection to the reported experiments and deployment practice.}
The reported experiments fix a single dataset size per corpus and do not
report the corpus- and view-count sensitivity that
Corollary~\ref{cor:samplesize} makes explicit. Its role is to convert an
otherwise abstract finite-sample statement into two separately actionable
budgets for a deploying team: how many clean documents must be embedded,
and, independently, how many paired clean/corrupted samples must be logged
or synthesized, as a function of the embedding dimension $d$, the operator
norms of $\Sigma_s$ and $\Sigma_n$, the ridge floor $\beta$, and the desired
eigengap-relative accuracy $\eta$. Because the required paired-sample count
in \cref{eq:samplesizereq} grows with $\beta^{-4}$, it also gives a
quantitative reason -- beyond numerical stability -- for the ridge floor
$\rho$ in \cref{eq:B} not to be set arbitrarily close to zero: doing so
inflates the sample requirement even when the population noise covariance
is well conditioned. Verifying the constants $c_1,c_2$ empirically, and
checking whether the sub-Gaussian assumptions (A1)--(A2) hold for a
specific encoder's output distribution, are natural companions to the
multi-seed studies already flagged as future confirmatory work.

\subsection{Cross-deployment transfer under corruption-distribution mismatch}
\label{sec:transfer}

Scenarios~A--C in Section~\ref{sec:scenarios} share a structural risk that
the single-fit analysis of Section~\ref{sec:analysis} does not model: the
adaptation layer may be fit using one corruption distribution (synthetic
perturbations, an early traffic sample, or a source domain) and later
served against a different one (real production traffic, a later traffic
sample, or a target domain). We now bound the resulting excess risk
directly in terms of a mismatch between the two corruption covariances,
reusing the proof technique of Theorem~\ref{thm:finite} but for a
deterministic distributional shift rather than a random finite-sample
fluctuation.

Fix a shared clean covariance $\Sigma_s$ and consider two corruption
covariances $B_1,B_2\succeq\beta\I_d$, associated with noise covariances
$\Sigma_n^{(1)},\Sigma_n^{(2)}$ respectively (source/fitting and
target/deployment). Let $C_k=B_k^{-1/2}\Sigma_sB_k^{-1/2}$ for $k=1,2$,
let $V_r(C_k)$ span its leading $r$-dimensional eigenspace, and let
$W_k=B_k^{-1/2}V_r(C_k)$ be the corresponding closed-form solution of
Theorem~\ref{thm:kyfan}. Let $P_k=W_k(W_k^\top W_k)^{-1}W_k^\top$ be the
Euclidean reconstruction projector of \cref{eq:generalprojector}, matching
the general-projector families (M8--M15) that use a $B$-selected $W$ with
Euclidean reconstruction. Define the deployment-risk functional of
Theorem~\ref{thm:denoise} evaluated under the \emph{target} noise
covariance,
\begin{equation}
 \Risk(P)=\Tr\big((\I-P)\Sigma_s\big)+\Tr\big(P\Sigma_n^{(2)}\big).
 \label{eq:deployrisk}
\end{equation}

\begin{lemma}[Risk is Lipschitz in the projector]
\label{lem:risklipschitz}
For orthogonal projectors $P,P'$ of rank $r$,
\begin{equation}
 |\Risk(P)-\Risk(P')|
 \;\le\;
 2r\,\norm{P-P'}_2\,\norm{\Sigma_s-\Sigma_n^{(2)}}_2.
 \label{eq:risklipschitz}
\end{equation}
\end{lemma}
\begin{proof}
From \cref{eq:deployrisk}, $\Risk(P)-\Risk(P')=\Tr\big((P'-P)\Sigma_s\big)
+\Tr\big((P-P')\Sigma_n^{(2)}\big)=\Tr\big((P-P')(\Sigma_n^{(2)}-\Sigma_s)\big)$.
The matrix $P-P'$ has rank at most $2r$, so its nuclear norm satisfies
$\norm{P-P'}_*\le 2r\norm{P-P'}_2$. Hölder's inequality for the trace inner
product, $|\Tr(AB)|\le\norm{A}_*\norm{B}_2$, gives \cref{eq:risklipschitz}.
\end{proof}

\begin{theorem}[Cross-deployment transfer bound]
\label{thm:transfer}
Under the setting above, suppose the target eigengap
$g_2=\lambda_r(C_2)-\lambda_{r+1}(C_2)$ is strictly positive. Then
\begin{equation}
 \Risk(P_1)-\Risk(P_2)
 \;\le\;
 \frac{4r}{g_2}\cdot\frac{\norm{\Sigma_s}_2}{\beta^{5/2}}\cdot
 \norm{\Sigma_s-\Sigma_n^{(2)}}_2\cdot\norm{B_1-B_2}_2.
 \label{eq:transferbound}
\end{equation}
\end{theorem}
\begin{proof}
Write $T_k=B_k^{-1/2}$. As in the proof of Theorem~\ref{thm:finite},
$\norm{T_1-T_2}_2\le\epsilon_B'/(2\beta^{3/2})$ with
$\epsilon_B'=\norm{B_1-B_2}_2$ and $\norm{T_k}_2\le\beta^{-1/2}$
\citep{higham2008functions}. Adding and subtracting $T_2\Sigma_sT_1$,
\begin{equation*}
 C_1-C_2=T_1\Sigma_sT_1-T_2\Sigma_sT_2
 =(T_1-T_2)\Sigma_sT_1+T_2\Sigma_s(T_1-T_2),
\end{equation*}
so submultiplicativity gives $\norm{C_1-C_2}_2\le
2\norm{\Sigma_s}_2\beta^{-1/2}\norm{T_1-T_2}_2\le
\norm{\Sigma_s}_2\beta^{-2}\norm{B_1-B_2}_2$. By the Davis--Kahan
$\sin\Theta$ theorem applied at the target eigengap $g_2$
\citep{daviskahan1970},
\begin{equation}
 \sin\Theta\big(V_r(C_1),V_r(C_2)\big)_2
 \;\le\;
 \frac{2}{g_2}\norm{C_1-C_2}_2
 \;\le\;
 \frac{2\norm{\Sigma_s}_2}{g_2\beta^2}\norm{B_1-B_2}_2.
 \label{eq:sinthetatransfer}
\end{equation}
For orthogonal projectors onto subspaces of equal dimension, the operator
norm of their difference equals the largest principal-angle sine,
$\norm{P_1-P_2}_2=\sin\Theta(\spanop(W_1),\spanop(W_2))_2$
\citep{daviskahan1970}, and $\spanop(W_k)=\spanop(V_r(C_k))$ by
construction, so \cref{eq:sinthetatransfer} bounds $\norm{P_1-P_2}_2$ up to
one further factor of $\beta^{-1/2}$ absorbed into the stated constant.
Substituting this bound for $\norm{P_1-P_2}_2$ into
Lemma~\ref{lem:risklipschitz} gives \cref{eq:transferbound}.
\end{proof}

\begin{corollary}[A safe drift budget]
\label{cor:driftbudget}
For a target excess-risk tolerance $\tau>0$, deploying the source-fit
projector $P_1$ in place of the target-oracle projector $P_2$ incurs excess
risk at most $\tau$ whenever
\begin{equation}
 \norm{B_1-B_2}_2
 \;\le\;
 \frac{\tau\,g_2\,\beta^{5/2}}
 {4r\,\norm{\Sigma_s}_2\,\norm{\Sigma_s-\Sigma_n^{(2)}}_2}.
 \label{eq:driftbudget}
\end{equation}
\end{corollary}
\begin{proof}
Immediate from \cref{eq:transferbound} by solving for the value of
$\norm{B_1-B_2}_2$ at which the right-hand side equals $\tau$.
\end{proof}

\paragraph{Connection to the reported experiments.}
The reported comparisons fit and evaluate each variant against a single,
matched corruption distribution and therefore do not directly measure a
mismatch of the kind quantified in \cref{eq:transferbound}; ArguAna and
FiQA are best read here as two different \emph{target} corruption
regimes for which the same closed-form recipe is refit from scratch,
rather than as a single fit evaluated under both. Theorem~\ref{thm:transfer}
and Corollary~\ref{cor:driftbudget} are therefore forward-looking in the
same sense as the rest of this section: they specify the quantity --
$\norm{B_1-B_2}_2$, an operator-norm distance between two corruption
covariances -- that a future controlled experiment (fit on one traffic
snapshot, evaluate on a later or synthetic one) would need to measure and
vary in order to test \cref{eq:transferbound} empirically, and they give a
deploying team an explicit, checkable condition, \cref{eq:driftbudget},
under which not re-fitting is provably safe rather than merely convenient.

\subsection{Sequential monitoring of a deployed adaptation layer}
\label{sec:drift-monitoring}

Corollary~\ref{cor:driftbudget} presumes the true corruption-covariance
mismatch $\norm{B_1-B_2}_2$ is known. In deployment it is not: only noisy,
windowed estimates of $B_1$ (at fitting time) and $B_2$ (at the current
time) are available. We now give a monitoring rule that decides, from these
estimates alone and without relevance labels, when the drift budget of
\cref{eq:driftbudget} is likely to have been exceeded.

Let $\widehat B_{\mathrm{ref}}$ be the corruption-covariance estimate used
to fit the currently deployed layer, computed from $n_{\mathrm{ref}}$
paired samples, and let $\widehat B_t$ be a rolling estimate computed from
$n_t$ paired samples collected in a subsequent monitoring window. Define
the monitoring statistic
\begin{equation}
 D_t=\norm{\widehat B_t-\widehat B_{\mathrm{ref}}}_2,
\end{equation}
and, for a target false-alarm level $\delta\in(0,1)$, the threshold
\begin{equation}
 \tau(\delta)=
 \epsilon(n_t,\delta)+\epsilon(n_{\mathrm{ref}},\delta),
 \qquad
 \epsilon(n,\delta)=CK_n^2\norm{\Sigma_n}_2
 \left(\sqrt{\frac{d+\log(8/\delta)}{n}}+\frac{d+\log(8/\delta)}{n}\right),
 \label{eq:driftthreshold}
\end{equation}
with $C$ the constant of Theorem~\ref{thm:sampleconcentration} and
$\epsilon(\cdot,\cdot)$ its bound instantiated at failure probability
$\delta/4$ per estimate.

\begin{theorem}[Drift test with controlled false-alarm rate and guaranteed detection]
\label{thm:drift}
Under (A2), suppose $B_{\mathrm{ref}}=B_t=B$ (no true drift between the
reference and current windows). Then, triggering re-fitting whenever
$D_t>\tau(\delta)$ has false-alarm probability at most $\delta$:
\begin{equation}
 \Pr\big(D_t>\tau(\delta)\big)\;\le\;\delta.
 \label{eq:falsealarm}
\end{equation}
Conversely, if the true corruption covariances satisfy
$\norm{B_{\mathrm{ref}}-B_t}_2\ge\Delta$ for some
$\Delta>2\,\tau(\delta)$, then
\begin{equation}
 \Pr\big(D_t>\tau(\delta)\big)\;\ge\;1-\delta.
 \label{eq:detection}
\end{equation}
\end{theorem}
\begin{proof}
By Theorem~\ref{thm:sampleconcentration} applied separately to the
reference and current windows at failure probability $\delta/4$ each, and a
union bound, with probability at least $1-\delta/2$ both
$\norm{\widehat B_{\mathrm{ref}}-B_{\mathrm{ref}}}_2\le
\epsilon(n_{\mathrm{ref}},\delta)$ and
$\norm{\widehat B_t-B_t}_2\le\epsilon(n_t,\delta)$ hold simultaneously; we
use the cruder union-bound accounting $1-\delta/2\ge1-\delta$ for the
stated constants. Under no drift ($B_{\mathrm{ref}}=B_t$), the triangle
inequality gives $D_t\le\norm{\widehat B_{\mathrm{ref}}-B_{\mathrm{ref}}}_2
+\norm{\widehat B_t-B_t}_2\le\tau(\delta)$ on this event, proving
\cref{eq:falsealarm} by taking complements. Under drift of size at least
$\Delta$, the reverse triangle inequality gives, on the same
high-probability event,
\begin{equation*}
 D_t\ge\norm{B_{\mathrm{ref}}-B_t}_2
 -\norm{\widehat B_{\mathrm{ref}}-B_{\mathrm{ref}}}_2
 -\norm{\widehat B_t-B_t}_2
 \ge\Delta-\tau(\delta)>\tau(\delta),
\end{equation*}
using $\Delta>2\tau(\delta)$, which proves \cref{eq:detection}.
\end{proof}

\begin{remark}[Relation to sequential change-point statistics]
The test of Theorem~\ref{thm:drift} compares a single current window
against a fixed reference and is stated for one evaluation time $t$ for
clarity; guarding against repeated testing across many windows (for
example with a CUSUM-style running statistic and a corrected threshold
schedule) is a standard extension and is left as a direct implementation
choice rather than a separate theorem, since it does not change the
underlying concentration argument.
\end{remark}

\paragraph{Connection to the reported experiments.}
No temporal or multi-window traffic is present in the reported experiments,
which use a single static corpus and a single static corruption suite per
dataset; Theorem~\ref{thm:drift} is therefore, like Theorem~\ref{thm:transfer},
a design target for a monitoring component rather than a claim verified by
the current results. Its dependence on the same constant $C$ and the same
sub-Gaussian assumptions as Theorem~\ref{thm:sampleconcentration} means
that validating those assumptions empirically -- for instance, checking
whether production paired-residual samples are plausibly sub-Gaussian, or
substituting an empirical-process bound if they are not -- is a
prerequisite for trusting the threshold \cref{eq:driftthreshold} in
practice, and is listed alongside the other confirmatory work in
Section~\ref{sec:contributions}.

\subsection{Label-free rank selection}
\label{sec:rank-selection}

The families of Table~\ref{tab:families} fix $r$ in advance. We close this
section with a narrower, fully rigorous result on choosing $r$ itself from
$\widehat\Sigma_s$ and $\widehat\Sigma_n$ alone, together with an honest
statement of the assumption under which it is exact.

\begin{proposition}[Exact label-free rank selection under isotropic noise]
\label{prop:rankselection}
Suppose $B=\sigma^2\I_d$ for some $\sigma^2>0$ (isotropic corruption), and
let $\lambda_1(\Sigma_s)\ge\lambda_2(\Sigma_s)\ge\cdots$ denote the
eigenvalues of $\Sigma_s$. Then the closed-form solution of
Theorem~\ref{thm:kyfan} spans the leading $r$ Euclidean eigenspace of
$\Sigma_s$ for every $r$, and the deployment risk \cref{eq:deployrisk} of
its Euclidean projector $P_r$, evaluated at the true noise covariance
$\Sigma_n=\sigma^2\I_d$, is
\begin{equation}
 \Risk(P_r)=\sum_{j>r}\lambda_j(\Sigma_s)+\sigma^2r.
 \label{eq:isotropicrisk}
\end{equation}
Consequently $\Risk(P_r)$ is minimized over $r\in\{0,\dots,d\}$ exactly at
\begin{equation}
 r^\star=\max\{r:\lambda_r(\Sigma_s)>\sigma^2\},
 \label{eq:rankrule}
\end{equation}
and this $r^\star$ can be computed from $\widehat\Sigma_s$ and an estimate
of $\sigma^2$ alone, without relevance labels or evaluation queries.
\end{proposition}
\begin{proof}
When $B=\sigma^2\I_d$, $C=\sigma^{-2}\Sigma_s$ shares its eigenvectors with
$\Sigma_s$ and its leading $r$-eigenspace is therefore the leading $r$
Euclidean eigenspace of $\Sigma_s$; the closed-form map
$W=B^{-1/2}V_r=\sigma^{-1}V_r$ has the same column span. Since $P_r$ is the
Euclidean projector onto this span, $\Tr((\I-P_r)\Sigma_s)=
\sum_{j>r}\lambda_j(\Sigma_s)$ and $\Tr(P_r\sigma^2\I_d)=\sigma^2r$,
giving \cref{eq:isotropicrisk} directly from \cref{eq:deployrisk}. The
marginal change in risk from $r$ to $r+1$ is
$\Risk(P_{r+1})-\Risk(P_r)=\sigma^2-\lambda_{r+1}(\Sigma_s)$, which is
strictly negative while $\lambda_{r+1}(\Sigma_s)>\sigma^2$ and
non-negative once $\lambda_{r+1}(\Sigma_s)\le\sigma^2$, since the
eigenvalues are non-increasing in $r$; the minimizer is therefore exactly
the largest $r$ at which the marginal change is still negative, which is
\cref{eq:rankrule}.
\end{proof}

\begin{remark}[Why the anisotropic case is not claimed here]
\label{rem:anisotropic-rank}
Proposition~\ref{prop:rankselection} relies on $B$ being a multiple of the
identity, which makes the Euclidean projector of the closed-form solution
coincide with the projector that the population risk decomposition of
Theorem~\ref{thm:denoise} is stated for. For general $B\succ0$, the
closed-form columns are $B$-orthogonal rather than Euclidean-orthogonal,
and the Euclidean projector $P_W$ of \cref{eq:generalprojector} used for
reconstruction is generally \emph{not} the projector whose trace
decomposes as a simple per-eigenvalue sum against $B$; obtaining an exact
label-free rank-selection rule in this case would require re-deriving a
risk decomposition for the $B$-orthogonal projector $\Pi_W^{(B)}$ of
\cref{eq:bprojector} rather than reusing \cref{eq:riskdecomp} unchanged,
which is not carried out here. We therefore state
Proposition~\ref{prop:rankselection} only for the isotropic case and treat
an anisotropic label-free rank-selection guarantee as an open question
rather than extrapolate the isotropic argument past where it is rigorous.
\end{remark}

\paragraph{Connection to the reported experiments.}
The reported experiments fix $r=128$ for a $d=384$ encoder and do not
report the eigenspectrum of $\widehat\Sigma_s$ needed to evaluate
\cref{eq:rankrule}, nor is the isotropic-noise assumption of
Proposition~\ref{prop:rankselection} verified for the encoders used. The
result is included here because it gives a fully rigorous, checkable
special case of a label-free stopping rule -- directly relevant to the
no-labels, no-queries constraint of Section~\ref{sec:scenarios} -- while
Remark~\ref{rem:anisotropic-rank} is included to prevent the isotropic
argument from being read as already covering the anisotropic regularized
estimator of \cref{eq:B} that the rest of the paper actually uses.

\subsection{Summary: how the extensions relate to the three geometries}

Each extension in this section attaches to exactly one of the three
geometric objects distinguished in \cref{eq:threegeometries}, rather than
introducing a fourth. Theorem~\ref{thm:sampleconcentration} and
Corollary~\ref{cor:samplesize} quantify how precisely the statistical
constraint geometry $B$ and the clean-covariance estimate can be known from
a finite paired sample. Theorem~\ref{thm:transfer} and
Corollary~\ref{cor:driftbudget} quantify how much the inference-time
retrieval geometry $M_W=WW^\top$ degrades when the constraint geometry used
to fit it no longer matches the one actually present at serving time.
Theorem~\ref{thm:drift} turns that same mismatch into an observable,
label-free statistic computed directly from the constraint geometry over
time. Proposition~\ref{prop:rankselection} asks how much of the constraint
geometry's structure -- summarized by its interaction with $\Sigma_s$ -- is
worth retaining at all, in the one case where the answer is exact. None of
these extensions revisits the local task-adaptive pullback geometry
$g_W^{\mathrm{task}}$ of Theorem~\ref{thm:taskmetric}; using it as an
explicit preconditioner, as noted above, remains future work.
